% Generated draft: requires mathematical and source review.
\documentclass[11pt]{article}
\usepackage[margin=25mm]{geometry}
\usepackage{fontspec}
\setmainfont{Latin Modern Roman}
\usepackage{amsmath,amssymb,mathtools,unicode-math}
\setmathfont{Latin Modern Math}
\usepackage{enumitem,needspace,xurl,hyperref}
\setlength{\parindent}{0pt}
\setlength{\parskip}{5pt}
\setlength{\emergencystretch}{3em}
\newcommand{\N}{\mathbb{N}}
\newcommand{\Z}{\mathbb{Z}}
\newcommand{\Q}{\mathbb{Q}}
\newcommand{\R}{\mathbb{R}}
\newcommand{\C}{\mathbb{C}}
\newcommand{\F}{\mathbb{F}}
\newcommand{\A}{\mathbb{A}}
\newcommand{\PP}{\mathbb{P}}
\newcommand{\E}{\mathbb{E}}
\newcommand{\eps}{\varepsilon}
\newcommand{\dd}{\,\mathrm{d}}
\begin{document}
\section*{UnsolvedMath ID: 30000536}
\textbf{Source problem code:} catalogue entry
\subsection{Definitions and mathematical statement}
The source asks how to measure, or assign a framed-bordism value to, the noncompact flow-line moduli space $M(P,Q)$ for nonconsecutive critical points in its Morse-complex construction. The supplied assessment interprets nonconsecutive in terms of indices; exact ambient and framing data are not supplied.

In the standard finite-dimensional Morse--Smale setting, let $f$ be Morse on a closed smooth manifold with a metric for which stable and unstable manifolds are transverse. The unparameterized moduli space is
\[
 M(P,Q)=\{\gamma:\R\to X:\dot\gamma=-\nabla f(\gamma),
 \ \lim_{t\to-\infty}\gamma(t)=P,
 \ \lim_{t\to+\infty}\gamma(t)=Q\}/\R,
\]
where $\R$ acts by time translation. If nonempty, its dimension is $r=\operatorname{ind}(P)-\operatorname{ind}(Q)-1$. The higher-index-gap problem has $r\geq1$, rather than the zero-dimensional signed counts used in the ordinary differential.

In this framework noncompactness can arise from broken trajectories. The standard compactification has boundary strata of the form
\[
 M(P,R_1)\times M(R_1,R_2)\times\cdots\times M(R_j,Q),
\]
with decreasing indices. A larger index gap alone does not force actual breaking; relevant intermediate trajectories must exist. Compactness and this corner structure depend on the stated Morse--Smale/compactness hypotheses and are not asserted for an unspecified general flow.

Ordinary framed bordism $\Omega_r^{\mathrm{fr}}$ uses compact closed stably framed $r$-manifolds. A noncompact $M(P,Q)$, or a compactification with nonempty boundary, does not automatically give such a class. One needs coherent framing and boundary/gluing data, a suitable relative bordism or stable-homotopy target, or another invariant which accounts for the broken strata. Orientations alone do not specify a framing. If ``measured'' instead means a numerical measure/invariant, its normalization, invariance and compatibility with gluing must likewise be specified.

The research request is to give this higher-moduli assignment in the report's intended Morse construction, not simply count an infinite set or choose an arbitrary volume. The literature triage checked 21 August 2026 calls it a formulation-level problem and supplies no precise target or general resolution. No independent current-status claim is made.
\subsection{Short English statement}
How can positive-dimensional Morse trajectory spaces, with ends corresponding to broken flows, contribute a meaningful measure or bordism value? Ordinary framed bordism requires compact closed framed manifolds, so boundary, framing and gluing data must be supplied. The source leaves the intended invariant unspecified.
\subsection{Sources}
\begin{itemize}
\item Local UnsolvedMath registry, permanent ID 30000536, exact source code catalogue entry; \emph{Bordism-Valued Moduli Spaces of Morse Flow Lines}. Proposed year 2006, status open; corrected\_verified editing preserves the higher-moduli measurement/bordism question.
\item Supplied report citation: \emph{Geometric Group Theory, Hyperbolic Dynamics and Symplectic Geometry (2007)}, Oberwolfach Reports, DOI 10.4171/owr/2006/33, pp. 1991--2058, \url{https://doi.org/10.4171/owr/2006/33}. Title year differs from DOI/proposed year.
\item Supplied literature triage checked 2026-08-21, citing only ``OWR Geometric Group Theory report'', \url{https://doi.org/10.4171/owr/2006/33}. It describes the nonconsecutive-index question as formulation-level in the reported construction. The dated OWR-derived curation requests independent verification of time-sensitive status and cited claims.
\end{itemize}
\end{document}
